\section{Motivation}

Reducing noise and vibration in lightweight structures is a common engineering challenge. Thin plates and panels are everywhere in cars, aircraft and machines, and they radiate sound when they vibrate. A lightweight way to reduce that radiated sound at a chosen frequency range is to attach small tuned mass resonators to the plate, the principle of the classical tuned vibration absorber \parencite{denhartog1956vibrations}. Each resonator is a small spring, damper and mass combination placed at a position on the plate, and a few of them placed at the right spots can noticeably change how much sound power a plate radiates. This kind of structure, a plain plate carrying a handful of local resonators, is a simple example of what is often called an acoustic or vibroacoustic metamaterial: its acoustic behaviour comes from the small attached resonators, not from the base material itself \parencite{liu2000locally}. Resonator-based structures of this kind are studied, for example, for broadband sound isolation in panels \parencite{lu2020membrane} and for acoustic liners in aero-engines \parencite{neubauer2023liner}.

Designing such a layout is not simple. The number of resonators, their masses, their tuning frequencies, their positions and their damping all interact with each other and with the plate's own vibration modes. An exact, closed-form solution exists only for a few simple cases, such as the simply supported rectangular plate used in this thesis; most realistic structures need a full multiphysics simulation, which is accurate but far too slow to try out many layouts by hand or by brute-force search. Classical design optimisation, for example topology optimisation \parencite{bendsoe2004topology}, automates the search, but it needs a forward evaluation for every candidate design it considers. Even where a closed-form solver is available, as it is here, it still only answers one direction of the problem: given a design, what is its acoustic behaviour? An engineer usually wants the opposite: given a target acoustic behaviour, what design achieves it? This second, harder direction is called generative design, or inverse design, and it is the central problem this thesis addresses.

The approach taken here is to replace the slow physics solver with fast, trained neural network surrogates, and then to build surrogates that work in both directions with one set of weights, so the same trained model can predict a spectrum from a design and propose a design from a target spectrum.

\section{This Thesis in Context of Research at the Chair}

Neural operators for generative design of tuned vibration absorbers are also being developed concurrently at the Chair of Vibroacoustics of Vehicles and Machines, as part of a project funded by the German Research Foundation (DFG) under the Priority Programme SPP 2353, ``Daring More Intelligence, Design Assistants in Mechanics and Dynamics'' \parencite{dogu2026gamm}. That project uses a physics-informed DeepONet to predict the full displacement field of a plate carrying one or several tuned vibration absorbers, with the motivation of reducing low-frequency noise from ducted fans in electric urban air mobility. Extending that model from a single absorber to several absorbers is reported as work in progress, and applying an invertible neural operator so the same model can also propose an absorber layout for a target displacement field is named as future work.

This thesis pursues a closely related goal, generative design of resonator layouts with invertible neural operators, but along a complementary route: it predicts the equivalent radiated power spectrum rather than the full displacement field, it does not add a physics-based loss term, and its resonator designs are drawn with Latin hypercube sampling across five variables per resonator rather than varying a single parameter at a time. The two efforts share the same long-term aim at the Chair: fast neural operators that support an engineer during the design of resonant structures, in both the forward and the inverse direction. The models developed in this thesis are not physics-grounded in the sense of solving the governing equations or using a physics-based loss term; they only use physics-aware input features, described in Chapter~\ref{ch:concept}, and are trained purely on data from the solver.

\section{Literature Review}

\subsection{Neural Operators for Forward Prediction}

A neural operator is a neural network trained to approximate a mapping between functions rather than a mapping between fixed-size vectors, so that once trained it can be evaluated on a new input function without retraining \parencite{kovachki2023neural}. The DeepONet architecture \parencite{lu2021deeponet} was one of the first popular neural operators: it splits the network into a branch that encodes the input and a trunk that encodes the output location, and combines the two with a dot product. The Fourier neural operator \parencite{li2021fno} instead learns a convolution directly in the Fourier domain, which works well for problems with a regular grid; a systematic comparison of the two on a range of problems is given by \textcite{lu2022comparison}. The wavelet neural operator \parencite{tripura2022wno} and the Laplace neural operator \parencite{cao2023lno} follow the same idea with a wavelet transform and a pole-residue (Laplace) representation instead, which can be a better fit for signals with localised features or transients, such as resonance peaks. The graph neural operator \parencite{li2020graphno} represents the input as a graph and learns a message-passing kernel \parencite{gilmer2017mpnn}, which fits naturally when the input is a set of discrete objects, such as the resonators in this thesis, rather than values on a grid. The Set Transformer \parencite{lee2019settransformer} offers another way to handle a set of unordered objects, using attention \parencite{vaswani2017attention} instead of message passing. Finally, SIREN \parencite{sitzmann2020siren} uses sine activations instead of the usual ReLU or tanh activations, which makes a network better at representing oscillatory signals and their derivatives. This thesis builds ten forward architectures that combine these ideas, adapted to take a set of resonators as input and to output the equivalent radiated power spectrum.

\subsection{Invertible and Generative Models}

A model that goes in both directions with one set of weights needs a building block that is exactly invertible. The basic idea can be seen with two numbers $x_1$ and $x_2$: instead of computing some function $y = f(x_1, x_2)$ that cannot generally be undone, compute a sum and a difference, $y_1 = x_1 + x_2$ and $y_2 = x_1 - x_2$. Knowing $y_1$ and $y_2$ is enough to recover $x_1$ and $x_2$ exactly, $x_1 = (y_1+y_2)/2$ and $x_2 = (y_1-y_2)/2$, so this simple transform is invertible by construction, even though it mixes $x_1$ and $x_2$ together. Normalizing flows build more expressive invertible transforms from the same principle. RealNVP \parencite{dinh2017realnvp} splits its input in two halves and transforms one half conditioned on the other with an affine map whose inverse can be written down directly. A variational autoencoder \parencite{kingma2014vae} is used in this thesis for a different purpose: not for exact invertibility, but to turn a single predicted design into a distribution of plausible designs, since many different resonator layouts can produce a similar spectrum. Invertible networks of this kind have also been used to analyse ill-posed inverse problems directly, because they can return a distribution of the parameters that explain an observation instead of a single guess \parencite{ardizzone2019inn}. Several recent papers combine these ideas with neural operators. \textcite{kaltenbach2023invertible} build an invertible DeepONet by making the branch network a RealNVP flow, so that the same weights map a design to its coefficients and back. \textcite{long2025invertiblefno} build an invertible Fourier neural operator with coupling blocks in the latent space and a variational autoencoder for the posterior over designs, trained in three stages: first the invertible core, then the autoencoder, then both together. More general approaches to inverse design with neural operators exist as well. Design-GenNO \parencite{zang2026designgenno} encodes microstructures into a low-dimensional latent space regularised by a normalising-flow prior and decodes them with neural operators, so that several design objectives can be addressed without retraining, and its physics-informed training reduces the need for labelled data. The inverse basis-to-basis framework $B2B^{-1}$ \parencite{thorpe2025generalizable} separates the representation of functions from the inverse map: it learns neural basis functions for the input and output spaces and then trains deterministic, invertible or probabilistic inverse models on the resulting coefficients, and it is compared against invertible neural operator baselines. A recent survey \parencite{nelsen2025operator} places such end-to-end inverse operator learning in the wider context of operator learning for inverse problems and discusses how noise in the data affects the inversion. The invertible architectures developed in this thesis follow the staged training recipe of the invertible Fourier neural operator.

\subsection{Sampling the Design Space}

Any data-driven surrogate is only as good as the dataset it is trained on. When a design has several variables, drawing them independently at random tends to leave gaps and clusters in the design space. Latin hypercube sampling \parencite{mckay1979lhs} avoids this by stratifying each variable into evenly spaced bins and drawing one sample per bin, which spreads the samples much more evenly across the whole design space than plain random sampling. This thesis uses Latin hypercube sampling for every design variable, and Chapter~\ref{ch:concept} shows a concrete example of the kind of bias that appears when a variable is not sampled directly but only derived from others.

\section{Objectives}

The objectives of this thesis are:

\begin{enumerate}
  \item Develop forward neural operators, a set of ten architectures (DON, DNO, DCO, SIREN, FNO, WNO, LNO, GNO, STO and NN) that predict the equivalent radiated power spectrum of a resonator design, so that they can serve as fast, accurate surrogates for the physics solver.
  \item Evaluate how accurately the forward surrogates reproduce the physics solver's spectra, and how their accuracy relates to their size and architecture.
  \item Develop invertible neural operators, iDON and an iFNO family (iDNO, iDCO, iFNO, iLNO, iGNO, iSTO), that use a single set of weights for both directions: predicting a spectrum from a design, and proposing a design for a target spectrum, enabling generative design of resonator layouts.
  \item Evaluate the invertible operators' ability to propose resonator designs for a target spectrum, and discuss how uncertainty quantification over manufacturing tolerances can check whether a proposed design stays close to its target once small, unavoidable manufacturing errors are taken into account.
\end{enumerate}

Objectives 1 and 2 produce and evaluate the forward surrogates, whose cores the invertible operators reuse. Objectives 3 and 4, generative design through invertible neural operators, are the main objective of this thesis.

\section{Thesis Outline}

Chapter~\ref{ch:theory} covers the physics and machine learning background needed to follow the rest of the thesis. Chapter~\ref{ch:concept} describes the design space, the dataset and the ten forward and seven invertible architectures used in this thesis. Chapter~\ref{ch:results} reports the forward surrogates' accuracy, for two and three resonators, and the invertible operators' generative-design performance. Chapter~\ref{ch:conclusions} summarises the findings and lists the next steps toward full generative design. A separate frequency-extrapolation experiment, which is not part of the main objective of this thesis, is reported in Appendix~A.
